% ============================================================
% chapter_06_conclusion.tex - Chapter 6, Conclusion
%
% Purpose: summarises the empirical findings, the methodological
%          finding (full reproducibility on free cloud infrastructure),
%          and the recommendations for the wider research community.
% Sections: 6.1 Summary of Findings, 6.2 Implications,
%           6.3 Recommendations.
% Included by main_thesis.tex.
% License: CC BY 4.0 (see ../LICENSE).
% ============================================================

\chapter{CONCLUSION}\label{chap:conclusion}

This thesis set out to present a hybrid bias correction framework
for daily satellite precipitation and to establish, dimension by
dimension, what it corrects and where it is bounded. It delivered
the result in two registers: what the framework achieves as a
scientific result, and what it provides as a reproducible
methodological artefact. This chapter summarises
the findings, states their implications, and sets out
recommendations for users, researchers, and the institutions that
maintain the underlying data.


% ============================================================
\section{Summary of Findings}\label{sec:conclusion-summary}
% ============================================================

The scientific finding is a stage-resolved account of what the
LSEQM+DL (Linear Scaling, Empirical Quantile Mapping, Generalized
Pareto Distribution tail correction, and Deep Learning refinement
through a Convolutional Neural Network) framework corrects and
what it leaves unmoved. Applied to the Integrated Multi-satellitE Retrievals for the
Global Precipitation Measurement (GPM) mission (IMERG) Late Run
(IMERG-L) over Indonesia across 2001 to 2025 and validated against
172 independent Indonesian Meteorological, Climatological, and
Geophysical Agency (BMKG) stations over 2001 to 2021, the pipeline moves three
of the four verification pillars cleanly toward the gauge target.

The value-adjustment
pillar reaches a relative bias near zero and a standard deviation
ratio of unity; the distribution-alignment pillar reaches a
Kolmogorov-Smirnov $p$-value at which the corrected and gauge
distributions are no longer statistically distinguishable; and
the extreme-preservation pillar reaches $Q_{99}$ within one
percent of the gauge target. The fourth pillar, event detection,
shows a designed trade-off: the aggregate probability of
detection falls as the dry-day-matching mechanism suppresses
spurious light rain. The threshold-stratified view, however, shows the
LSEQM+DL pipeline overtaking Linear Scaling at and above the
$20$~mm/day very-heavy-precipitation threshold of the
internationally standardised extremes indices.
Each stage moves the metrics matched to its design dimension and
leaves the others alone.

The most consequential result is the one metric that
does not move. Measured per pixel against the CPC-UNI calibration
target, the Pearson correlation between paired daily values stays
near $0.35$ across all three corrected stages, a shift smaller
than the single-pixel sampling uncertainty. The same flat pattern
holds against the independent BMKG stations, station-by-station in
the Taylor analysis and month-by-month across the seasonal cycle,
and it survives the correction of the calendar-window offset
(Section~\ref{sec:results-intercomparison}). This is not a failure of the implementation. It is the
predicted behaviour of any marginal correction family: rank
preservation and variance inflation without time-dependent skill
together leave the daily correlation set by the day-by-day
pairing quality of the raw satellite retrieval, which lies
upstream of every step in the pipeline.

The corrected
in-sample value of $0.348$ against CPC-UNI is of the order reported
for daily IMERG over the Indonesian Maritime Continent
\citep{ref:ramadhan2022ts}, and well below the median near $0.75$
the same product reaches against the far denser Chinese gauge
analysis \citep{ref:tang2020}; the calendar-window diagnostic
accounts for most of that difference. The framework's reach is
therefore precise: a marginal-correction framework
can take a daily field all the way to the gauge marginal
distribution at every percentile, and no distance at all along
the axis of day-by-day timing.

These results resolve the three hypotheses of
Section~\ref{sec:intro-hypotheses}. H1 is supported: each stage
moves its matched pillar toward the gauge target, with the
Convolutional Neural Network acting as a small, spatially targeted
refinement of the extreme tail
(Section~\ref{sec:results-stages}). H2 is supported against both
references and at both date labels: the daily correlation does not
improve beyond the sampling uncertainty across the corrected
stages, against the in-sample CPC-UNI target and the independent
BMKG stations alike. H3 holds in the GPM era: the low daily
correlation against BMKG is substantially a calendar-label
artifact, recoverable by matching the satellite to the gauge day,
and the recovery is specific to the era in which the retrieval
resolves the diurnal cycle.

The methodological finding is that this entire account is
independently reproducible. The pipeline is released as a
configuration-driven set of source modules and numbered
notebooks that run unchanged in a local environment or on a
free-tier cloud instance. The correction-to-visualisation sequence
was re-executed
on a free Google Colab instance for the Bali subdomain across all
36 dekadal windows in 72.1 minutes on a standard Central
Processing Unit (CPU) runtime. That establishes that an independent
reader can reconstruct the results for a self-contained region
without institutional compute access; the full Indonesian domain
scales up with pixel count and is correspondingly longer. The two findings are of equal
weight: the scientific result is credible because it is
reproducible, and the reproducible artefact is worth releasing
because it carries a result worth re-verifying.


% ============================================================
\section{Implications}\label{sec:conclusion-implications}
% ============================================================

For operational use, the central implication is that the
corrected product should be matched to the metrics that moved.
The corrected daily IMERG-L is suitable for the broad class of
applications that depend on distributional properties:
dekadal-to-monthly water budgets, drought indicators such as the
Standardized Precipitation Index and its evapotranspiration
variant, extreme-event frequency analysis, threshold-exceedance
climatology, and hydrological model calibration against multi-year
statistics. It is not suitable, on its own, for applications that
depend on the precise calendar day of an event: real-time
flood-event nowcasting, single-event forecasting, and day-1
hydrological initialisation inherit the timing ceiling of the raw
retrieval and must combine the corrected product with sub-daily
disaggregation, move to ensemble methods, or accept the timing
limitation. Correcting at the daily scale remains necessary so
that the dekadal, monthly, and seasonal aggregates derived from
the daily field are bias-free.

For the open-science contribution, the implication is that the
verification posture argued for in this thesis can be applied by
others. The four-pillar framework with Pearson correlation
reported as a separate temporal-skill diagnostic, rather than
folded into a single skill score, lets a structural limit and a
methodological improvement appear in the same report. An
independent group can apply the same verification scaffolding to
a different product, a different reference, or a different region,
and can reproduce or challenge the present result on its own
terms.


% ============================================================
\section{Recommendations}\label{sec:conclusion-recommendations}
% ============================================================

\medskip\noindent\textbf{For users of corrected satellite
precipitation.}
Read the verification report before choosing an application.
A product that matches the gauge distribution is the right input
for distributional questions and the wrong input for
day-specific timing questions, and the two cases cannot be told
apart from an aggregate skill score alone. Where the application
is timing-sensitive and operates over the Maritime Continent, do
not reject the product on its daily correlation before checking
the calendar: the UTC-aligned daily window is a known contributor
to the timing ceiling, and the recovery available from a
locally-aligned re-aggregation, quantified in
Section~\ref{sec:discussion-subdaily}, costs nothing elsewhere in
the pipeline.

\medskip\noindent\textbf{For researchers extending the method.}
Start with the calendar window. Of the four extensions set out in
Section~\ref{sec:discussion-paths}, re-aggregating the daily
archive to the local-observation window is the only one that
leaves the pipeline untouched, and it addresses the larger of the
two components of the timing ceiling. The remaining three, which
reach the residual retrieval-skill component, all require methods
outside the marginal correction family and are correspondingly
more costly. Three further extensions are worth stating as
recommendations rather than as findings: a joint search over the
three sensitivity parameters with a strict temporal hold-out,
which this study explored only one parameter at a time; a
transfer experiment to a second tropical region, which would test
a portability the framework's design allows but does not
establish; and the fully-convolutional Stage~4 head described in
Section~\ref{sec:methods-cnn}, which matched the dense model on
the Bali subdomain at a fraction of the parameter count and is
the natural replacement at full resolution.

\medskip\noindent\textbf{For BMKG, CPC, and Institut Pertanian Bogor (IPB).}
The reproducibility contribution depends on continued access to
the underlying archives and on the maintenance of the released
code against dependency drift. A containerised distribution and a
versioned data snapshot would make the reproducibility claim
durable beyond the dependency state recorded at release. More
broadly, the gap this thesis addresses, namely the rarity of
released code in the bias correction literature, is one that
institutional data-sharing and reproducibility standards are well
placed to close, and the framework released here is offered as a
concrete example of what such standards can produce.
